% ============================================================================
% LEHRERHINWEISE: parecer y conocer
% ============================================================================
% Abgeleitet von: LE-01-parecer-conocer.tex
% Enthält: Kurzplan, Differenzierung [B]/[S]/[E], Lösungen, Dialogkarten
% ============================================================================

\documentclass[11pt,a4paper]{article}
\usepackage[utf8]{inputenc}
\usepackage[T1]{fontenc}
\usepackage[ngerman]{babel}
\usepackage{geometry}
\usepackage{fancyhdr}
\usepackage{xcolor}
\usepackage{tcolorbox}
\usepackage{enumitem}
\usepackage{tabularx}
\usepackage{longtable}
\usepackage{array}

% Layout
\geometry{left=2cm, right=2cm, top=2cm, bottom=2cm}
\definecolor{fach}{HTML}{c41e3a}
\definecolor{lightgray}{HTML}{f5f5f5}
\definecolor{basisgreen}{HTML}{2a7a4b}
\definecolor{standardblue}{HTML}{2c5aa0}
\definecolor{erweiterungpurple}{HTML}{7b1fa2}
\definecolor{loesung}{HTML}{c62828}

\tcbuselibrary{skins,breakable}

% Kopfzeile
\pagestyle{fancy}
\fancyhf{}
\fancyhead[L]{\textbf{Spanisch Klasse 9} -- LH-01}
\fancyhead[R]{\textcolor{fach}{\textbf{NUR FÜR LEHRKRÄFTE}}}
\renewcommand{\headrulewidth}{0.4pt}

% Befehle
\newcommand{\basis}[1]{\textcolor{basisgreen}{\textbf{[B]} #1}}
\newcommand{\standard}[1]{\textcolor{standardblue}{\textbf{[S]} #1}}
\newcommand{\erweiterung}[1]{\textcolor{erweiterungpurple}{\textbf{[E]} #1}}

\begin{document}

% ============================================================================
% TITEL
% ============================================================================
\begin{center}
{\Large\textcolor{fach}{\textbf{Lehrerhinweise: parecer y conocer}}}\\[0.3cm]
{\small Qué pasa 2, S. 18--21 | Doppelstunde 90 Min.}
\end{center}

\vspace{0.5cm}

% ============================================================================
% KURZPLAN
% ============================================================================
\section*{Kurzplan (90 Min.)}

\begin{longtable}{|p{1cm}|p{2cm}|p{6cm}|p{2cm}|p{2cm}|}
\hline
\textbf{Min.} & \textbf{Phase} & \textbf{Aktivität} & \textbf{SF} & \textbf{Material} \\
\hline
\endhead

\multicolumn{5}{|l|}{\textbf{1. Stunde: Texterschließung + parecer}} \\
\hline
0--5 & Einstieg & Bildimpuls Patchwork-Familie & PL & Bild \\
\hline
5--20 & Erarbeitung & Text T2 (S. 18) lesen, Aufg. 10a & EA/GA & Buch \\
\hline
20--30 & Grammatik & parecer im Text finden, Regel & PL & Tafel \\
\hline
30--40 & Übung & Partnerdialog mit parecer & PA & Dialogkarten \\
\hline
40--45 & Sicherung & Merkkasten diktieren & PL & Tafel \\
\hline

\multicolumn{5}{|l|}{\textbf{2. Stunde: conocer + Anwendung}} \\
\hline
0--5 & Wiederholung & Quiz parecer & PL & -- \\
\hline
5--15 & Grammatik & conocer: yo-Form, + a & PL & Tafel \\
\hline
15--30 & Übung & AB Aufgaben 1--4 & EA & AB \\
\hline
30--40 & Anwendung & Partnerdialog (beide Verben) & PA & -- \\
\hline
40--45 & Abschluss & Zusammenfassung, HA & PL & Tafel \\
\hline
\end{longtable}

% ============================================================================
% DIFFERENZIERUNG
% ============================================================================
\section*{Differenzierung}

\begin{tcolorbox}[colback=green!5, colframe=basisgreen, title=\textbf{[B] Basis}]
\begin{itemize}
    \item Konjugationstabelle (parecer + conocer) als Hilfe
    \item Dialogkarten mit Lücken statt freier Produktion
    \item Bei Aufg. 4: Satzanfänge vorgeben
    \item Mehr Zeit für Textarbeit, Vorentlastung Wortschatz
\end{itemize}
\end{tcolorbox}

\begin{tcolorbox}[colback=blue!5, colframe=standardblue, title=\textbf{[S] Standard}]
\begin{itemize}
    \item Selbstständige Bearbeitung mit Satzanfängen
    \item Partnerdialog mit vorgegebenen Themen
    \item Bei Aufg. 4: Beispiele als Orientierung
\end{itemize}
\end{tcolorbox}

\begin{tcolorbox}[colback=purple!5, colframe=erweiterungpurple, title=\textbf{[E] Erweiterung}]
\begin{itemize}
    \item Freie Dialoge ohne Vorgaben
    \item Regel conocer + a selbst erklären
    \item Zusatz: 5 Sätze über eigene Familie (HA)
    \item Jugendsprache aus S. 20 verwenden
\end{itemize}
\end{tcolorbox}

% ============================================================================
% ERWARTETE FEHLER
% ============================================================================
\section*{Erwartete Fehler und Reaktionen}

\begin{tabular}{|p{5cm}|p{8cm}|}
\hline
\textbf{Fehler} & \textbf{Reaktion} \\
\hline
*Me parece los deberes difíciles & ``parecer'' braucht Singular/Plural! $\rightarrow$ Me parece\textbf{n} (Plural) \\
\hline
*Conoco (statt conozco) & c $\rightarrow$ zc vor o! (wie traducir $\rightarrow$ traduzco) \\
\hline
*Conozco María (ohne a) & Merkregel: ``Personas con \textbf{a}!'' -- Person-Symbol an Tafel \\
\hline
parecer/gustar verwechselt & ``Me gusta = ich mag'' vs. ``Me parece = ich finde'' \\
\hline
\end{tabular}

% ============================================================================
% LÖSUNGEN AB
% ============================================================================
\section*{Lösungen Arbeitsblatt}

\textbf{Aufgabe 1:} (6 P.)
\begin{enumerate}[label=\alph)}]
    \item parecen -- parecen
    \item parece -- parece
    \item parece
    \item parece -- parece
\end{enumerate}

\textbf{Aufgabe 2:} (7 P.)
\begin{enumerate}[label=\alph)]
    \item conozco
    \item conoces
    \item conocemos
    \item conocen
    \item conoce
    \item conocéis
    \item conozco
\end{enumerate}

\textbf{Aufgabe 3:} (6 P.)
\begin{enumerate}[label=\alph)]
    \item a (Person)
    \item -- (Sache)
    \item a (Personen)
    \item -- (Sache)
    \item a (Person)
    \item -- (Sache)
\end{enumerate}

\textbf{Aufgabe 4:} (7 P.) -- Individuelle Lösungen, z.B.:
\begin{itemize}
    \item Me parece interesante. / Me parecen difíciles.
    \item Conozco a mi profesor. / Conozco la ciudad.
\end{itemize}

% ============================================================================
% DIALOGKARTEN (Kopiervorlage)
% ============================================================================
\newpage
\section*{Dialogkarten (Kopiervorlage)}

\begin{center}
\fbox{\parbox{6cm}{
\textbf{Karte A: parecer}

Pregunta: ¿Qué te parece...?
\begin{itemize}
    \item Gonzalo
    \item la clase de español
    \item los deberes
\end{itemize}

Responde con: Me parece(n) + Adjektiv
}}
\hspace{1cm}
\fbox{\parbox{6cm}{
\textbf{Karte B: parecer}

Pregunta: ¿Qué te parece...?
\begin{itemize}
    \item Clara
    \item el fútbol
    \item las vacaciones
\end{itemize}

Responde con: Me parece(n) + Adjektiv
}}
\end{center}

\vspace{1cm}

\begin{center}
\fbox{\parbox{6cm}{
\textbf{Karte C: conocer}

Pregunta: ¿Conoces a...?
\begin{itemize}
    \item mi hermano
    \item los amigos de Clara
    \item Nicolás
\end{itemize}

Responde: Sí, conozco a... / No, no conozco a...
}}
\hspace{1cm}
\fbox{\parbox{6cm}{
\textbf{Karte D: conocer}

Pregunta: ¿Conoces...?
\begin{itemize}
    \item la calle Trinidad
    \item esta canción
    \item el centro de Cádiz
\end{itemize}

Responde: Sí, conozco... / No, no conozco...
}}
\end{center}

% ============================================================================
% TAFELANSCHRIEB
% ============================================================================
\section*{Tafelanschrieb}

\begin{center}
\fbox{\parbox{14cm}{
\textbf{\large parecer y conocer}

\vspace{0.3cm}

\textbf{1. parecer} (wie gustar!)

me/te/le/nos/os/les + parece (Sg.) / parecen (Pl.)

\textit{¿Qué te parece Gonzalo?} -- \textit{Me parece simpático.}

\vspace{0.3cm}

\textbf{2. conocer} (kennen)

conozco -- conoces -- conoce -- conocemos -- conocéis -- conocen

\vspace{0.2cm}

\textbf{conocer + a} bei \textbf{PERSONEN}!

¿Conoces \textbf{a} María? -- Sí, conozco \textbf{a} María.

¿Conoces la ciudad? -- Sí, conozco la ciudad.
}}
\end{center}

\end{document}
